\Askhsh{13752}{4}
\begin{ylh}{1}
    \item 3.2. 1ο Κριτήριο ισότητας τριγώνων
    \item 3.12. Τριγωνική ανισότητα
    \item 4.4. Γωνίες με πλευρές παράλληλες.
\end{ylh}
% ===================================================================
%  Αρχείο: 13752.tex (Fragment)
%  Αυτόματη μετατροπή μέσω doc2latex.py
% ===================================================================

ΘΕΜΑ 4

Σε τρίγωνο ΑΒΓ με $\widehat{Β}$ \textless{} 90\textsuperscript{ο} θεωρούμε τυχαίο σημείο Δ της πλευράς ΑΓ. Φέρουμε τμήμα ΔΕ ίσο και παράλληλο με την πλευρά ΒΓ και από το σημείο Ε φέρουμε τμήμα ΕΖ ίσο και παράλληλο με την πλευρά ΑΒ, όπως φαίνεται στο σχήμα.

\begin{alist}
\item Ένας μαθητής κάνει τους παρακάτω διαδοχικούς συλλογισμούς. Να χαρακτηρίσετε Σ (Σωστό) ή Λ (Λάθος) κάθε έναν από αυτούς.

\begin{alist}
\item
  Οι γωνίες Δ$\widehat{Ε}$Ζ και Α$\widehat{Β}$Γ είναι γωνίες με πλευρές παράλληλες.
\item
  Οπότε Δ$\widehat{Ε}$Ζ = Α$\widehat{Β}$Γ.
\item
  Τα τρίγωνα ΔΕΖ και ΑΒΓ είναι ίσα.
\item
  Το τμήμα ΔΖ είναι ίσο με το τμήμα ΑΓ.

\monades{08}

\item Να αιτιολογήσετε τους χαρακτηρισμούς σας (Σ ή Λ) που αφορούν τους ισχυρισμούς 2. και 3. \monades{10}

\item Αν στα δεδομένα παραλείψουμε τη συνθήκη $\widehat{Β}$ \textless{} 90\textsuperscript{ο}, να συγκρίνετε τα τμήματα ΑΓ και ΔΖ για τα διάφορα είδη της γωνίας $\widehat{Β}$ και να δικαιολογήσετε τις απαντήσεις σας. \monades{07}


\begin{center}
\begin{tikzpicture}[scale=1]
\tkzDefPoint(0,4){A}
\tkzDefPoint(-2,0){B}
\tkzDefPoint(2,0){C}
\tkzDefPointOnLine[pos=0.4](A,C) \tkzGetPoint{D}
\tkzDefPointBy[translation=from B to C](D) \tkzGetPoint{E}
\tkzDefPointBy[translation=from B to A](E) \tkzGetPoint{Z}

\draw[l3] (A)--(B)--(C)--cycle;
\draw[l3] (D)--(E);
\draw[l3] (Z)--(E);
\draw[l3] (D)--(Z);

\tkzDrawPoints(A,B,C,D,E,Z)
\tkzLabelPoint[above](A){$A$}
\tkzLabelPoint[left](B){$B$}
\tkzLabelPoint[below](C){$\Gamma$}
\tkzLabelPoint[above right](D){$\Delta$}
\tkzLabelPoint[right](E){$E$}
\tkzLabelPoint[above](Z){$Z$}

\tkzMarkSegments[mark=s|](B,C D,E)
\tkzMarkSegments[mark=s||](A,B Z,E)
\end{tikzpicture}
\end{center}
\end{alist}
\end{alist}
